\documentclass[11pt]{article}

\usepackage[a4paper,margin=1in]{geometry}
\usepackage{amsmath,amssymb}
\usepackage{hyperref}

\hypersetup{
    colorlinks=true,
    linkcolor=blue,
    citecolor=blue,
    urlcolor=blue
}

\title{Derivation of the Wasserstein Gradient Flow Equation}
\author{}
\date{}

\begin{document}

\maketitle

In Euclidean space, the gradient flow of a function $f$ is
\[
\dot x_t=-\nabla f(x_t).
\]
This says that the point moves in the direction of steepest decrease, measured
using the Euclidean inner product.

In Wasserstein space, points are probability densities $\rho_t$. A curve of
densities moves through a velocity field $v_t$ satisfying the continuity
equation:
\[
\partial_t\rho_t+\operatorname{div}(\rho_t v_t)=0.
\]
The Wasserstein metric measures the size of the velocity by kinetic energy:
\[
\|v_t\|_{\rho_t}^2
=
\int_{\mathbb R^d}|v_t(x)|^2\rho_t(x)\,dx.
\]

\section*{Derivative of the Functional}

Let $\mathcal F$ be a functional on densities. Along a curve $\rho_t$,
\[
\frac{d}{dt}\mathcal F(\rho_t)
=
\int_{\mathbb R^d}
\frac{\delta\mathcal F}{\delta\rho}(\rho_t)(x)
\partial_t\rho_t(x)\,dx.
\]
Using the continuity equation,
\[
\frac{d}{dt}\mathcal F(\rho_t)
=
-\int_{\mathbb R^d}
\frac{\delta\mathcal F}{\delta\rho}(\rho_t)
\operatorname{div}(\rho_t v_t)\,dx.
\]
Integrating by parts gives
\[
\frac{d}{dt}\mathcal F(\rho_t)
=
\int_{\mathbb R^d}
\nabla
\left[
\frac{\delta\mathcal F}{\delta\rho}(\rho_t)
\right]
\cdot v_t\,\rho_t\,dx.
\]

\section*{Identifying the Wasserstein Gradient}

The Wasserstein inner product of two velocity fields is
\[
\langle u,v\rangle_{\rho_t}
=
\int_{\mathbb R^d}u\cdot v\,\rho_t\,dx.
\]
The previous identity shows that the velocity-field representative of the
Wasserstein gradient of $\mathcal F$ is
\[
\nabla
\left[
\frac{\delta\mathcal F}{\delta\rho}(\rho_t)
\right].
\]
Therefore the negative-gradient velocity is
\[
v_t
=
-
\nabla
\left[
\frac{\delta\mathcal F}{\delta\rho}(\rho_t)
\right].
\]

\section*{Plugging Into the Continuity Equation}

Substituting this velocity into
\[
\partial_t\rho_t+\operatorname{div}(\rho_t v_t)=0
\]
gives
\[
\partial_t\rho_t
=
\operatorname{div}
\left(
\rho_t
\nabla
\left[
\frac{\delta\mathcal F}{\delta\rho}(\rho_t)
\right]
\right).
\]
This is the Wasserstein gradient flow equation.

\end{document}
